\documentclass[11pt]{article}
\usepackage{microtype,amsmath,amssymb,amsthm,mathtools,booktabs,array}
\usepackage[margin=1.05in]{geometry}
\usepackage[colorlinks=true,linkcolor=blue!45!black,citecolor=blue!45!black,urlcolor=blue!45!black]{hyperref}
\usepackage{xcolor}
\hypersetup{pdftitle={The dimension-degree boundary for Gaussian moment vanishing},pdfsubject={Working mathematical manuscript}}
\newtheorem{theorem}{Theorem}[section]
\newtheorem{proposition}[theorem]{Proposition}
\newtheorem{lemma}[theorem]{Lemma}
\newtheorem{corollary}[theorem]{Corollary}
\theoremstyle{definition}\newtheorem{definition}[theorem]{Definition}
\theoremstyle{remark}\newtheorem{remark}[theorem]{Remark}
\newcommand{\E}{\mathbb E}
\newcommand{\C}{\mathbb C}
\newcommand{\tr}{\operatorname{tr}}
\newcommand{\wt}{\operatorname{wt}}
\newcommand{\G}{\mathsf G}
\DeclareMathOperator{\Span}{span}
\title{Formal stationary balls and a leading-coefficient obstruction}
\author{}\date{17 November 2026: retained written proof}
\begin{document}\maketitle
This is a written proof over a valued field, including arbitrary
nonconvergent formal Laurent coefficients. It is not an all-ring
statement. It gives an all-degree one-variable exponential-module
criterion and two cubic Gaussian applications. Priority is unassessed;
independent replay, machine-formal verification and external review
are not claimed. No finite computation proves its completeness lemmas.

\section{Gaussian conventions for the applications}
Let $T$ be standard Gaussian and $Z,W$ circular coordinates with
$\E[Z^iW^j]=\delta_{ij}i!$, independent of $T$. For a polynomial
$P$, set $\Phi_Q=\sum_{m\geq0}\E[QP^m]t^m/m!$. Purity means
$\Phi_1=1$; safety means every $\Phi_Q$ is a polynomial.
These are formal series, without an integral-convergence hypothesis.
Gaussian integration by parts is $\E[(\partial_T-T)H]=0$.

For any polynomial $h(z,T)$, let $g=th(g,T)$ be the unique small
formal solution. The exact circular identity is
\[
 \E_{Z,W}[A(Z,T)e^{tWh(Z,T)}]
       =\frac{A(g,T)}{1-th_Z(g,T)}.
\]
Indeed the left side is $\sum_m t^m[z^m]A(z,T)h(z,T)^m$.
It is the formal residue of $A(z,T)/(z-th(z,T))$, whose unique
small pole is $g$. This also proves the identity for a
$t$-dependent polynomial-coefficient formal numerator $A$.

\section{The full formal frame and its Gaussian applications}
\subsection{A formal leading-coefficient criterion}\label{sec:stationaryballs}

For a polynomial phase, the remaining shrinking argument can be
replaced by a valuation calculation. The construction below includes
critical points that collide, remain multiple, or escape to infinity.
The key is to group them before making a Gaussian approximation.

Write $v$ for the valuation of a formal parameter $r$, normalized by
$v(r)=1$. The valuation of a polynomial is the least valuation of its
coefficients. We use the formal Newton--Puiseux theorem: finitely many
algebraic elements over $\C((r))$ belong to one field
$\C((r^{1/N}))$. Convergence is not required~\cite[Chapter~2]{Wall}.

\begin{lemma}[Polynomial separation of formal balls]\label{lem:ballseparation}
Let finitely many disjoint balls have centers $c_B$ in a Puiseux field
and radii specified by $v(x-c_B)\geq\rho_B$. After one common
ramification, put $\delta_B=r^{\rho_B}$. For any one ball $B$ and
integer $j\geq0$, there is a polynomial $q_{B,j}(x)$ over the Laurent
field whose substitution $x=c_C+\delta_C y$ has integral coefficients
for every ball $C$, and whose reduction is $y^j$ on $B$ and zero on
every other ball.
\end{lemma}
\begin{proof}
The finitely many valuations $v(c_B-c_C)$ give a finite tree of
clusters of centers. At a node with scale $r^d$, the coordinate
$u=(x-c)/r^d$ has distinct constant reductions $a_1,\ldots,a_s$
on its children. The Lagrange polynomial
\[
 e_i(u)=\prod_{k\ne i}\frac{u-a_k}{a_i-a_k}
\]
reduces to one on child $i$ and zero on every other child.

Start at the chosen leaf with $((x-c_B)/\delta_B)^j$. Induct upward.
Suppose a polynomial $h$ already has the required reductions on the
chosen child. On each of the other child balls its coefficients have
some finite lower valuation. The restriction of $e_i$ there has
strictly positive valuation. Consequently $e_i^M h$, for a sufficiently
large fixed integer $M$, has positive coefficient valuation on all
those children. On the chosen child $e_i^M=1+$ positive-order terms,
so the existing reductions are unchanged. There are finitely many
nodes, and this constructs the required polynomial. A node with just
one leaf requires no filter. This also explains why independently
normalized Lagrange polynomials, with no control of their growth on
other scales, would not suffice.
\end{proof}

\begin{theorem}[Formal stationary balls]\label{thm:stationaryballs}
Let $\varphi(r,x)\in\C((r))[x]$ have degree $D\geq2$, and let
$h(r)\in\C((r))\setminus\{0\}$. Put
\[
 M=\C((r))[x]/(\partial_x+\varphi_x)\C((r))[x],\qquad
 \nabla[q]=[q_r+(\varphi_r+h'/h)q].
\]
After a finite ramification, this module has a full frame of rank-one
formal exponential functionals. If $a_D$ is the leading coefficient
of $\varphi$, every zero-exponential functional in this frame takes
the seed $[1]$ to a series of valuation at least
\begin{equation}\label{eq:valuationbound}
 v(h)-\frac{v(a_D)}D.
\end{equation}
In particular, if $Dv(h)>v(a_D)$, the differential submodule generated
by $[1]$ has no horizontal quotient taking $[1]$ to the constant one.
\end{theorem}
\begin{proof}
The two twisted derivatives commute, so the parameter connection
is well defined on the quotient. A nonzero polynomial of degree $m$
has spatial twisted derivative of degree $m+D-1$. Reduction of
successive leading terms therefore gives a basis
$1,x,\ldots,x^{D-2}$ and dimension $D-1$.

Split $\varphi_x$, retaining root multiplicities. For a critical root
$\xi$, expand
\[
 \varphi(\xi+u)=c_\xi+\sum_{j=2}^D a_{\xi,j}u^j,
 \qquad
 \rho_\xi=\max_{2\leq j\leq D,\ a_{\xi,j}\ne0}
                    \frac{-v(a_{\xi,j})}{j}.
\]
Enlarge the ramification so that all these scales are integral, and
set $\delta_\xi=r^{\rho_\xi}$. The polynomial
\[
 \varphi(\xi+\delta_\xi y)-c_\xi
\]
has integral coefficients, no linear term, and at least one unit
coefficient. Its reduction $F_\xi(y)$ has degree $d_\xi\geq2$.
Here and below we may replace $r$ by a root of $r$ and renormalize
the valuation; every displayed strict inequality is unchanged.
As $a_{\xi,D}=a_D$,
\begin{equation}\label{eq:balllower}
 \rho_\xi\geq-v(a_D)/D.
\end{equation}

\emph{The balls partition all critical roots.}
If another critical root $\zeta$ satisfies
$v(\zeta-\xi)\geq\rho_\xi$, translate the integral scaled polynomial
by $(\zeta-\xi)/\delta_\xi$. All its coefficients remain integral,
and the coefficient of degree $d_\xi$ remains a unit: higher-degree
terms have positive valuation. The linear term vanishes because
$\zeta$ is critical. It follows from the definition of the maximum
that $\rho_\zeta=\rho_\xi$. Thus their balls are identical.
If two of the balls intersect, one contains the other; the same
argument makes them identical. Otherwise they are disjoint.

The scaled derivative has reduction $F_\xi'$ of degree $d_\xi-1$.
It therefore has exactly $d_\xi-1$ integral roots, with multiplicity.
Here is a direct check of that count: factor it in a Puiseux field.
An integral root contributes a linear factor to the reduction;
for a root of negative valuation, factor out that root and the
remaining factor reduces to a constant. Multiplicativity of the
coefficient valuation leaves a nonzero constant scalar. Hence the
degree of the reduction counts exactly the integral roots.
The ball $B$ thus contains $d_B-1$ critical roots. Since
$\deg\varphi_x=D-1$,
\begin{equation}\label{eq:ballrankcount}
 \sum_B(d_B-1)=D-1.
\end{equation}

\emph{A functional for every dimension in a ball.}
The fixed polynomial $F_B$ has twisted quotient
\[
 H_B=\C[y]/(\partial_y+F_B')\C[y]
\]
with basis $1,y,\ldots,y^{d_B-2}$. Let $\ell^0_{B,j}$ be its dual
basis. Write
\[
 \varphi(\xi_B+\delta_B y)=c_B+F_B(y)+E_B(r,y),
 \qquad E_B\in r\C[[r]][y].
\]
For a polynomial amplitude $q$ define
\begin{equation}\label{eq:ballfunctional}
 \ell_{B,j}(q)=h\delta_B e^{c_B}
   \ell^0_{B,j}\!\left(e^{E_B(r,y)}q(\xi_B+\delta_B y)\right).
\end{equation}
Each coefficient uses finitely many polynomial terms. The exponential
in this formula is purely formal. The product rule conjugates the
spatial twisted derivative to $\partial_y+F_B'$, so the functional
annihilates the defining relations of $M$.

We spell out the parameter check because it retains the moving volume.
At fixed $y$, put $w=\xi_B'+\delta_B'y$. Differentiating
\eqref{eq:ballfunctional}, then subtracting its value on
$q_r+(\varphi_r+h'/h)q$, leaves, inside the fixed functional,
\[
 e^{E_B}\left(wq_x+\frac{\delta_B'}{\delta_B}q+\varphi_xwq\right)
 =\frac1{\delta_B}(\partial_y+F_B')
                      \left(e^{E_B}wq\right).
\]
It is zero. Thus the functionals are horizontal, including the
Jacobian $\delta_B$ and normalization $h$.

\emph{Completeness, not just a rank count.}
Choose the polynomials $q_{B,j}$ of
Lemma~\ref{lem:ballseparation}, for $0\leq j<d_B-1$.
Divide each evaluation column by its scalar $h\delta_Be^{c_B}$.
The resulting matrix reduces to the identity: $E_B$ has positive
order and the reductions of the amplitudes are the chosen monomials
on one ball and zero elsewhere. The matrix is therefore invertible.
Leading-degree reduction gives $\dim M=D-1$, so
\eqref{eq:ballrankcount} accounts for the entire module.

Remove from $c_B$ its finite negative-power part $q_B(r)$.
The remaining exponential is an invertible formal power series,
up to a nonzero constant. Formula~\eqref{eq:ballfunctional} and the
invertible matrix give the asserted direct sum of rank-one factors
$e^{q_B}$ times powers of $r$. If $q_B=0$, its value on the seed
has valuation at least $v(h)+\rho_B$, proving
\eqref{eq:valuationbound}.

Finally, after ramification this is a semisimple differential module.
A quotient on the submodule generated by the seed extends to the
whole module. Nonzero exponential factors have no horizontal map
to the Laurent field, because the logarithmic derivative of a
Laurent series cannot be the derivative of a nonzero negative-power
exponent. On the zero part, every horizontal functional is a
constant combination of those in the full frame. If the bound is
strictly positive, their seed values cannot sum to one.
\end{proof}

The construction is a formal stationary-phase argument. Its classical
background includes Newton--Puiseux and the formal stationary-phase
theory of exponential connections; Sabbah's local Fourier--Laplace
formula is a primary reference for the latter~\cite{SabbahStationary}.
We do not identify that formula's scope with the particular proof
above or claim priority for this formulation. The proof here also
allows nonconvergent formal Laurent coefficients. It uses a valued
field, not an arbitrary coefficient ring.

\begin{corollary}[Leading-coefficient obstruction]\label{cor:leadingobstruction}
For a rational polynomial phase $\varphi(t,x)$ of degree $D\geq2$
and nonzero rational scalar normalization $h(t)$, let $v_\infty$
mean order at $\varepsilon=1/t=0$. If
\begin{equation}\label{eq:leadingcriterion}
 v_\infty(h^D/a_D)>0,
\end{equation}
its seed submodule cannot have a rational horizontal quotient with
seed value one. In particular any horizontal formal functional at
another point cannot take that seed to one.
\end{corollary}
\begin{proof}
Apply the theorem at infinity. A rational horizontal quotient remains
a quotient on extension to the formal Laurent field. A formal
functional with seed value one restricts to such a rational quotient:
its values on successive seed derivatives are $1,0,0,\ldots$.
\end{proof}

The inequality must be strict. For example
$\varphi=-x^2/(2t^2)$ and $h=1/t$ have
$v_\infty(h^2/a_2)=0$, and their normalized Gaussian period is a
nonzero constant. Thus the earlier adverse scaling example is also
an exact boundary test of the criterion.

\subsection{Closing the full rank-one $WT^2$ family}\label{sec:rankoneformal}

\begin{proposition}\label{prop:rankoneformal}
For $f\ne0$ and $\deg G\leq3$, the polynomial
\[
 P=W(a+bZ+cT+fT^2)+G(Z,T)
\]
is pure if and only if $b=0$ and $G=0$. Its pure members obey
$m>\deg Q$.
\end{proposition}
\begin{proof}
Put $q=1-bt$ and $A(T)=a+cT+fT^2$. Exact circular contraction
gives
\[
 J(p)=\frac1q\E_T\!\left[p(T)
             \exp\{tG(tA/q,T)\}\right],\qquad J(1)=\Phi_1.
\]
Gaussian integration by parts says that $J$ annihilates
$\partial_T+\varphi_T$; direct parameter differentiation makes it
horizontal for the polynomial phase and normalization
\[
 \varphi=tG(tA/q,T)-T^2/2,\qquad h=1/q.
\]
The phase has degree at least two, since it equals $-T^2/2$ at $t=0$.

Suppose first $b=0$ and write $G=\sum_{j\geq0}Z^jG_j(T)$.
If a term $A^jG_j$ attains the actual degree $D\geq2$ of the phase,
the leading coefficient is a polynomial in $t$ containing a
nonzero positive power. Distinct $j$ give distinct powers $t^{j+1}$,
so they cannot cancel identically. Its order at infinity is negative,
whereas $v_\infty(h)=0$. Criterion~\eqref{eq:leadingcriterion}
excludes purity. Because $\deg A=2$, every nonzero $G_j$ with
$j\geq1$ attains some degree at least two. The only remaining
possibility is $G=\alpha T+\beta$. Its exact unit series is
$\exp(\beta t+\alpha^2t^2/2)$, which is one only if $G=0$.

Suppose now $b\ne0$. Divide $G$ by $S=bZ+A(T)$ as a polynomial
in $T$:
\[
 G=R_0(Z)+TR_1(Z)+S(\alpha Z+\beta T+\gamma).
\]
The remainder has basis $1,T,Z,ZT,Z^2,Z^2T,Z^3$. At $Z=-A/b$
these monomials have distinct leading degrees $0,1,2,3,4,5,6$.
Thus if $V(T)=R(-A/b,T)$ has degree $n$, then
$\deg R(sA,T)\leq n$ for every scalar $s$; a constant $V$
forces $R$ to be that same constant. Write $g=tA/q$. The phase is
exactly
\[
 tR(g,T)+\alpha g^2+\beta gT+\gamma g-T^2/2.
\]
Since $g=-A/b+O(t^{-1})$, its first term is $tV+O(1)$, with
all these remainders of degree at most $n$.

If $\alpha\ne0$, the limiting remaining phase has degree four
with leading coefficient $\alpha f^2/b^2$. If $\alpha=0$,
$\beta\ne0$, it has degree three with leading coefficient
$-\beta f/b$. Otherwise its quadratic coefficient tends to
$-\gamma f/b-1/2$. A leading term from $tV$ has valuation $-1$;
a higher-degree nonzero limiting term has valuation zero.
The only flattened case has
$\alpha=\beta=0$, $\gamma=-b/(2f)$. Its exact quadratic
coefficient outside $tR$ is $-1/(2q)$, of valuation one.
Thus in every case the actual leading coefficient $a_D$ has
$v_\infty(a_D)\leq1$. This also covers $V$ constant or linear.
But $v_\infty(h)=1$ and $D\geq2$, so
$Dv_\infty(h)-v_\infty(a_D)\geq D-1>0$.
Purity is impossible by the same criterion.

Conversely $W A(T)$ has strictly negative circular weight.
For a monomial $Z^iW^jT^k$, a mixed moment can be nonzero only
when the number of $W$ factors balances the number of $Z$ factors;
in particular it vanishes for $m>\deg Q$.
\end{proof}

\subsection{A mixed chart with a complete formal proof}\label{sec:mixedformal}

One of the previously analytic exclusions admits a direct formal proof.
It is the whole family in the next proposition, not a generic choice
of its coefficients.

\begin{proposition}\label{prop:mixedformal}
For every $a\in\C$ and every $G\in\C[Z,T]$ of total degree at most
three, the polynomial
\[
 P=W(a+T+Z^2)+G(Z,T)
\]
has a nonzero positive pure Gaussian moment.
\end{proposition}
\begin{proof}
We construct the exact formal seed functional and apply
Corollary~\ref{cor:leadingobstruction}. No analytic integral is needed.

\emph{The formal source.}
The circular fixed point is $g=t(a+T+g^2)$, with determinant
$D=1-2tg$. Make the formal Gaussian change
\[
 T=Y-t^2(a+Y)^2,\qquad z=t(a+Y).
\]
It has $g=z$ and $dT/dY=D$. For any change $T(Y)=Y+O(t)$ with
polynomial formal coefficients, the identity
\[
 \E_T[F(T)]=\E_Y\!\left[F(T(Y))T'(Y)
             e^{(Y^2-T(Y)^2)/2}\right]
\]
holds coefficientwise. To prove it, use the homotopy
$T_u=Y+u(T-Y)$: its derivative is the Gaussian expectation of
$(\partial_Y-Y)\{F(T_u)\dot T_u e^{(Y^2-T_u^2)/2}\}$, which is zero.
Thus the determinant cancels and gives
\[
 \Phi_1=\E_Y[e^\Psi],\qquad
 \Psi=tG(z,Y-z^2)+Yz^2-z^4/2.
\]

Define
\[
 F=z/t-a-z^2,\qquad \varphi(t,z)=tG(z,F)-F^2/2,
\]
and let
\[
 M=K[z]/(\partial_z+\varphi_z)K[z],\qquad
 \nabla[q]=[q_t+(\varphi_t-1/t)q],\qquad K=\C(t).
\]
These are vector-space quotients. The formal map
$J(q)=\E_Y[q(t(a+Y))e^\Psi]$ annihilates
$(\partial_z+\varphi_z)K[z]$ by Gaussian integration by parts.
Differentiating at fixed $Y$ differs from applying $\nabla$ by
$J((\partial_z+\varphi_z)(zq/t))=0$.
So $J$ is horizontal and $J(1)=\Phi_1$. The term $-1/t$ is the
normalization $dz/t=dY$; removing it changes the problem.
Leading-degree reduction gives $\dim_K M=\deg_z\varphi-1$.

\emph{The exhaustive degree split.}
Put
\[
 V(z)=G(z,-a-z^2),\quad L(z)=G_T(z,-a-z^2),\quad
 B_1=z^2G_{TT}(z,-a-z^2)/2+z(a+z^2).
\]
If $s=[T^3]G$, direct Taylor expansion gives
\begin{equation}\label{eq:mixedphi}
 \varphi=tV+B_0+B_1/t+B_2/t^2,
 \quad B_0=zL-(a+z^2)^2/2,\quad B_2=sz^3-z^2/2.
\end{equation}
To check every degree drop, write $v_1=[ZT^2]G$,
$v_0=[T^2]G$ and $u_2=[Z^2T]G$. Then
\[
 [z^6]V=-s,\qquad [z^5]V=v_1,\qquad
 [z^4]V=v_0-u_2-3as.
\]
If $n=\deg V\geq4$, the degree of $\varphi$ is $n$ and each $B_j$
has degree at most $n$. If $n\leq3$, or $V$ is constant, these three
identities force $s=v_1=0$ and $u_2=v_0$. Hence
$B_0=-z^4/2+O(z^3)$, $\deg B_1\leq3$ and $B_2=-z^2/2$.
In that case $\deg_z\varphi=4$. There is no omitted higher-degree
term tending to zero at infinity.

If $n\geq4$, the leading coefficient of $\varphi$ has a nonzero
$tV_n$ term and therefore has order $-1$ at infinity. If $n\leq3$
or $V$ is constant, its leading coefficient is exactly $-1/2$ and
has order zero. In both cases $h=1/t$ has order one, so
\[
 v_\infty(h^D/a_D)>0.
\]
The leading-coefficient obstruction excludes $J(1)=1$. The degree
calculation was exhaustive, so this proves the proposition even at
multiple critical points and degree drops. The full formal frame,
including every escaping dimension, is supplied by
Theorem~\ref{thm:stationaryballs} rather than a new case-by-case
saddle calculation.
\end{proof}

\begin{thebibliography}{9}
\bibitem{SabbahStationary}
C.~Sabbah, \emph{An explicit stationary phase formula for the local
formal Fourier--Laplace transform}, in \emph{Singularities I},
Contemporary Mathematics \textbf{474}, American Mathematical Society,
2008, 309--330;
\href{https://arxiv.org/abs/0706.3570v2}{arXiv:0706.3570v2}.

\bibitem{Wall}
C.~T.~C.~Wall, \emph{Singular Points of Plane Curves}, Cambridge
University Press, 2004, Chapter~2, 15--38,
\href{https://doi.org/10.1017/CBO9780511617560.003}{doi:10.1017/CBO9780511617560.003}.

\end{thebibliography}
\end{document}
